\section{Annotation and Methods}

\subsection{Representation}

We use UD-style token-level representation in CoNLL-U, with per-token fields for form, lemma, UPOS, FEATS, HEAD, and DEPREL. Pipeline-side constraints enforce shape consistency before records are admitted to the frozen dataset.

\subsection{Operational robustness}

The annotation pipeline is designed for long-running, interruption-prone generation:
\begin{itemize}[leftmargin=1.25em]
  \item offset-based batch progression with resumable state,
  \item explicit retry queue for failed or malformed outputs,
  \item conservative parse fallbacks for near-valid JSON,
  \item cost and throughput monitoring for practical budget control.
\end{itemize}

\subsection{Model families}

We evaluate five model families:
\begin{enumerate}[leftmargin=1.25em]
  \item feature-based logistic parser with learned head labels,
  \item transformer parser with multilingual encoder (XLM-R),
  \item mBERT variants under matched split/evaluation settings,
  \item Stanza pretrained Greek baseline,
  \item custom-trained Stanza POS + depparse.
\end{enumerate}

Stanza-based components and comparisons use the standard toolkit implementation \cite{qi2020stanza}.
